% SPDX-FileCopyrightText: Copyright (c) 2016-2026 Objectionary.com
% SPDX-License-Identifier: MIT

\section{Lexical Layer}\label{sec:lexical}

EO source is read line by line.
The lexer decides character encoding, validates indentation,
  and turns each non-blank, non-comment line
  into a sequence of tokens that the syntactic layer (\cref{sec:syntax}) consumes.
This section specifies that surface: encoding, indentation, the token catalogue,
  the literal sub-languages (numeric, string, text), and the multi-line bytes form.

\subsection{Encoding and Line Endings}\label{sec:lexical:encoding}

Source text is a stream of \nospell{UTF-8} bytes.
Line terminators are either \ff{LF} (\ff{U+000A}) or
  \ff{CR LF} (\ff{U+000D U+000A});
  both are normalised to \ff{LF} before further processing.
Other line-terminator forms (bare \ff{CR}, \ff{NEL}, \ff{LS}, \ff{PS}) are rejected.

\subsection{Indentation}\label{sec:lexical:indent}

EO uses significant whitespace.
Indentation determines nesting:
  a line whose leading-space count is greater than the previous non-blank line's
  declares a child of that line's object.
Three rules govern indentation:

\begin{enumerate}
  \item The indent unit is \emph{two spaces}.
    The leading-space count of any non-blank line must be even;
    odd counts are rejected.
  \item Between two consecutive non-blank lines, the indent may increase
    by \emph{at most one level} (two spaces).
    A jump of two or more levels is rejected.
  \item The indent may decrease by any amount in one step.
    Decreasing closes one or more enclosing blocks at once.
  \item Tab characters in the leading whitespace are rejected.
    Tabs elsewhere on a line are permitted only inside string or text literals
    via the \ff{\textbackslash t} escape (\cref{sec:lexical:strings}).
\end{enumerate}

A legal sequence of indents:
\begin{ffcode}
[] > foo
  a > x          // indent 2 (level 1), parent foo at indent 0
    b            // indent 4 (level 2), one level deeper than parent
      c          // indent 6 (level 3)
  d > y          // decrease from 6 to 2, allowed
\end{ffcode}

An illegal sequence:
\begin{ffcode}
[] > foo
  a > x
      b          // rejected: indent jumped from 2 to 6 (skipping 4)
\end{ffcode}

Comment lines (those beginning with \ff{\#}) also obey the even-indent constraint;
  blank lines carry no indent of their own and never trigger a level change
  (cross-line consequences are spelled out in \cref{sec:wellformedness}).

\subsection{Tokens}\label{sec:lexical:tokens}

\Cref{tab:tokens} lists every lexical class.
Names of identifier-like tokens are written in \ff{UPPERCASE}
  to distinguish them from the syntactic non-terminals of \cref{sec:syntax}.

\begin{table*}
\capt{Lexical tokens of EO.}
\label{tab:tokens}
\small
\begin{tabularx}{\textwidth}{l X}
\toprule
Token & Description \\
\midrule
\ff{META}        & \ff{+}\,\ff{NAME} followed by zero or more space-separated parts;
                    each part is one or more non-whitespace characters,
                    and parts may contain \ff{:}, \ff{.}, \ff{-}, \ff{/}
                    (e.g.\ \ff{+rt jvm a.b.c:lib:1.0.0}). \\
\ff{COMMENTARY}  & \ff{\#} followed by the rest of the line.
                    Trailing whitespace is excluded at the token level. \\
\ff{NAME}        & A lowercase Latin letter (\ff{[a-z]}) followed by characters
                    other than space, line break, tab, and the punctuation set
                    \ff{,}, \ff{.}, \ff{|}, \ff{'}, \ff{:}, \ff{;}, \ff{!}, \ff{?},
                    \ff{]}, \ff{[}, \ff{\}}, \ff{\{}, \ff{)}, \ff{(},
                    and the cactus codepoint \ff{U+1F335} (\(\mathtt{\Yright}\)).
                    The cactus exclusion is structural: it is reserved for
                    parser-generated auto-names and cannot collide with user input. \\
\ff{PHI}         & The single character \ff{@}.
                    Maps to \(\varphi\) in the \(\varphi\)-calculus translation
                    (\cref{sec:semantics}). \\
\ff{RHO}         & The single character \ff{\textasciicircum{}}.
                    Maps to \(\rho\). \\
\ff{ROOT}        & The single character \ff{Q}.
                    Maps to \(\Phi\). \\
\ff{XI}          & The single character \ff{\$}.
                    Maps to \(\xi\). \\
\ff{INT}         & A signed integer literal (\cref{sec:lexical:numeric}). \\
\ff{FLOAT}       & A signed floating-point literal (\cref{sec:lexical:numeric}). \\
\ff{HEX}         & A hexadecimal integer literal (\cref{sec:lexical:numeric}). \\
\ff{STRING}      & A double-quoted string literal (\cref{sec:lexical:strings}). \\
\ff{TEXT}        & A triple-quoted text block (\cref{sec:lexical:strings}). \\
\ff{BYTES}       & A byte-sequence literal (\cref{sec:lexical:bytes}). \\
\ff{STAR}        & The single character \ff{*}. \\
\ff{ARROW}       & The single character \ff{>}. \\
\ff{DOT}         & The single character \ff{.}. \\
\ff{COLON}       & The single character \ff{:}. \\
\ff{SLASH}       & The single character \ff{/}. \\
\ff{CONST}       & The single character \ff{!}. \\
\ff{LSQ},\,\ff{RSQ} & The brackets \ff{[} and \ff{]}. \\
\ff{LB},\,\ff{RB}   & The parentheses \ff{(} and \ff{)}. \\
\bottomrule
\end{tabularx}
\end{table*}

The four single-character tokens \ff{PHI}, \ff{RHO}, \ff{ROOT}, and \ff{XI}
  are \emph{locator} tokens.
They denote, respectively,
  the special \(\varphi\) attribute of the enclosing formation (decoratee),
  the parent \(\rho\),
  the root namespace \(\Phi\), and
  the current scope \(\xi\) (self).
The semantic role of each locator is detailed in \cref{sec:semantics}.

Indentation is not a separate token: it is read directly from a line's leading
  whitespace by the syntactic layer.
There is no \ff{INDENT}/\ff{DEDENT} pair in the lexer.

\subsection{Numeric Literals}\label{sec:lexical:numeric}

Three numeric forms are recognised.

\paragraph{Integer.}
An \ff{INT} is an optional sign (\ff{+} or \ff{-}) followed by either
  the single digit \ff{0} or a digit in \ff{[1-9]} followed by any number of
  digits in \ff{[0-9]}.
Leading zeros on multi-digit literals are forbidden: \ff{07}, \ff{007},
  \ff{+07}, and \ff{-07} are all lexical errors.

\paragraph{Floating-point.}
A \ff{FLOAT} is an optional sign, one or more digits, a \ff{.},
  one or more digits, and an optional exponent \ff{(e|E)(+|-)?digits}.

\paragraph{Hexadecimal.}
A \ff{HEX} is the literal prefix \ff{0x} (lowercase) followed by one or more
  hexadecimal digits in \ff{[0-9a-fA-F]}.

\subsection{String and Text Literals}\label{sec:lexical:strings}

A \ff{STRING} is delimited by \ff{"} and may contain any character except an
  unescaped \ff{"}, backslash, carriage return, or newline,
  together with the escape sequences enumerated in \cref{tab:escapes}.
A \ff{STRING} is a single-line construct: the lexer rejects a literal that
  reaches end-of-line before the closing quote.

A \ff{TEXT} is delimited by \ff{"""} markers spanning multiple lines.
The opening \ff{"""} must appear with no other content on its line;
  body lines must be indented at least as deep as the opening;
  the closing \ff{"""} must appear on its own line at the opening indent.
The body text is the body lines joined by newline, with the opening indent
  stripped from each.
Body characters may be any character except an unescaped backslash,
  plus the escape sequences in \cref{tab:escapes}.

An unclosed text block is reported at end-of-stream as a recoverable error.

\begin{table*}
\capt{Escape sequences recognised in \ff{STRING} and \ff{TEXT} literals.}
\label{tab:escapes}
\small
\begin{tabularx}{\textwidth}{l X}
\toprule
Sequence & Meaning \\
\midrule
\ff{\textbackslash b}      & \ff{U+0008} (backspace) \\
\ff{\textbackslash t}      & \ff{U+0009} (horizontal tab) \\
\ff{\textbackslash n}      & \ff{U+000A} (line feed) \\
\ff{\textbackslash f}      & \ff{U+000C} (form feed) \\
\ff{\textbackslash r}      & \ff{U+000D} (carriage return) \\
\ff{\textbackslash "}      & \ff{"} \\
\ff{\textbackslash '}      & \ff{'} \\
\ff{\textbackslash \textbackslash} & \ff{\textbackslash} \\
\ff{\textbackslash NNN}    & Octal byte. Pattern: a backslash followed by an
                              optional digit in \ff{[0-3]}, an optional digit
                              in \ff{[0-7]}, and a required digit in \ff{[0-7]}.
                              The decoded value must not exceed \ff{0o377}
                              (255 decimal). \\
\ff{\textbackslash uXXXX}  & A Unicode codepoint as four hexadecimal digits.
                              For legacy compatibility, a sequence of one or
                              more \ff{u} characters between the backslash
                              and the four digits is accepted on input; the
                              canonical form has exactly one \ff{u}. \\
\bottomrule
\end{tabularx}
\end{table*}

Any other backslash sequence is a lexical error.

\subsection{Bytes Literals}\label{sec:lexical:bytes}

A \ff{BYTES} literal denotes a sequence of zero or more 8-bit bytes,
  written as pairs of hexadecimal digits separated by dashes.
Three forms exist:

\begin{itemize}
  \item \ff{-{}-} --- the empty byte sequence.
  \item \ff{BB-} --- a single byte (two hexadecimal digits) followed by a trailing dash.
    The trailing dash distinguishes a one-byte literal from an \ff{INT} or \ff{HEX}.
  \item \ff{BB-BB(-BB)*} --- two or more bytes joined by dashes.
\end{itemize}

The multi-byte form may continue across lines.
If a multi-byte chunk ends a line with a trailing dash, the literal is incomplete
  and continues on the next line with another \ff{BB(-BB)*} chunk.
Continuation may repeat any number of times.
The continuation indent of every chunk after the first must be at least as deep
  as the indent of the line that began the literal;
  a chunk at shallower indent terminates the literal and is a lexical error.
A comment line, a blank line, or any non-bytes content interleaved between
  continuation chunks is a lexical error.

Example of a multi-line bytes literal:
\begin{ffcode}
size.
  CA-FE-      // first chunk, ends with a trailing dash
  BE-BE       // continuation chunk
\end{ffcode}

The two indented lines under \ff{size.} form a single \ff{BYTES} literal
  \ff{CA-FE-BE-BE}, occupying one argument slot of the reversed dispatch
  \ff{size.} (the latter construct is defined in \cref{sec:syntax}).
A multi-line literal is one expression at one indent: the continuation
  affects only the lexer, not the indent-driven syntactic layer.
